% ============================================================
%  第二卷 重要符号（全书正文前材料，不编号）
%  底本：高教社中文版第8版
%  说明：
%  - OCR（full_merged.md 行 141–226）对该节漏行严重（三维表缺
%    dV,d𝒇,d𝒍、哈密顿函数等条目；四维表的表首说明行全部丢失），
%    本表按原书第 iv–v 页页面图（原书转换/pages/p-008.png、
%    p-009.png）逐行补全校排，行序忠实原书。
%  - 表内页码引用原书页码，现已改写为本重排本节号（§6/§23/§29/§32），避免页码失效。
% ============================================================

\chapter*{重要符号}
\phantomsection
\addcontentsline{toc}{chapter}{重要符号}

\section*{三维空间中的量}

\noindent 三维张量的指标用希腊字母表示\par
\medskip
\noindent
\begin{tabular}{@{}p{5.2cm}@{\qquad}p{9.6cm}@{}}
$\mathrm{d}V,\ \mathrm{d}\boldsymbol{f},\ \mathrm{d}\boldsymbol{l}$ & 体积元，面元，线元 \\
$\boldsymbol{p}$ 和 $\mathcal{E}$ & 粒子的动量和能量 \\
$\mathcal{H}$ & 哈密顿函数 \\
$\varphi$ 和 $\boldsymbol{A}$ & 电磁场的标势和矢势 \\
$\boldsymbol{E}$ 和 $\boldsymbol{H}$ & 电场强度和磁场强度 \\
$\rho$ 和 $\boldsymbol{j}$ & 电荷密度和电流密度 \\
$\boldsymbol{d}$ & 电偶极矩 \\
$\boldsymbol{m}$ & 磁偶极矩 \\
\end{tabular}

\section*{四维空间中的量}

\noindent 四维张量的指标用拉丁字母表示\par
\noindent 字母 $i,k,l,\cdots$ 的取值范围是 $0,1,2,3$\par
\noindent 采用带有号差 $(+---)$ 的度规\par
\noindent 指标的上移和下移法则——见 §6\par
\noindent 四维矢量的分量以 $A^{i}=(A^{0},\boldsymbol{A})$ 的形式列出\par
\medskip
\noindent
\begin{tabular}{@{}p{5.6cm}@{\qquad}p{9.2cm}@{}}
$e^{iklm}$ & 四阶反对称单位赝张量，并且 $e^{0123}=1$（定义见 §6） \\
$\mathrm{d}\Omega=\mathrm{d}x^{0}\mathrm{d}x^{1}\mathrm{d}x^{2}\mathrm{d}x^{3}$ & 四维体积元 \\
$\mathrm{d}S^{i}$ & 超曲面元（定义见 §6） \\
$x^{i}=(ct,\boldsymbol{r})$ & 四维径矢 \\
$u^{i}=\mathrm{d}x^{i}/\mathrm{d}s$ & 四维速度 \\
$p^{i}=(\mathcal{E}/c,\boldsymbol{p})$ & 四维动量 \\
$j^{i}=(c\rho,\rho\boldsymbol{v})$ & 四维电流 \\
$A^{i}=(\varphi,\boldsymbol{A})$ & 电磁场的四维势 \\
$F_{ik}=\dfrac{\partial A_{k}}{\partial x^{i}}-\dfrac{\partial A_{i}}{\partial x^{k}}$ &
  四维电磁场张量（分量 $F_{ik}$ 与 $\boldsymbol{E}$ 和 $\boldsymbol{H}$ 的分量之间的关系见 §23） \\
$T^{ik}$ & 四维能量动量张量（其分量的定义见 §29、§32） \\
\end{tabular}

\medskip
在引用本教程其他各卷的章节和公式时，卷号与书名的对应关系为：

第一卷：《力学》，俄文第五版，中文第一版；

第三卷：《量子力学（非相对论理论）》，俄文第六版，中文第一版；

第五卷：《统计物理学Ⅰ》，俄文第五版，中文第一版；

第六卷：《流体力学》，俄文第五版，中文第一版；

第八卷：《连续介质电动力学》，俄文第四版，中文第一版.
